% ============================================================
%  Поглавље 8: Закључак
% ============================================================
\section{Закључак}
\label{sec:zakljucak}

% --- Резиме доприноса рада ---
% Сумирај главне тачке из сваког поглавља у 2--3 реченице.

У овом раду разматрали смо НП-тешке проблеме кроз призму проблема
MaxCut и показали два комплементарна приступа: прецизно решавање
путем SMT решавача Z3, и апроксимативно решавање путем квантног
алгоритма QAOA.

Кључни закључци:

\begin{itemize}[leftmargin=1.5em]
  \item НП-тешки проблеми представљају суштинску препреку класичном
        рачунарству, а MaxCut је репрезентативни пример са богатом
        структуром применљивом на квантне технике.
  \item Z3 и SMT решавачи пружају прецизна решења за мале инстанце
        и служе као референца за верификацију квантних резултата.
  \item QAOA је природни квантни алгоритам за QUBO/Изингове проблеме;
        његова перформанса расте са дубином $p$, али га NISQ ограничења
        ограничавају на умерене дубине.
  \item Квантно рачунарство и аутоматско резоновање нису раздвојене
        области: хибридне петље, Гроверово убрзање и формална
        верификација квантних програма чине их природним саговорницима.
\end{itemize}

\subsection*{Отворена питања}

% --- Правци будућег истраживања ---
% Набројати 3--5 конкретних отворених проблема и праваца истраживања.

\begin{itemize}[leftmargin=1.5em]
  \item \textbf{Дубљи QAOA:} Да ли QAOA са великим $p$ може да
        надмаши Гоеманс-Вилијамсон апроксимацију ($0.878$) на
        свим класама графова?
  \item \textbf{Грешке и корекција:} Кад ће квантни хардвер са
        корекцијом грешака омогућити поуздани QAOA на великим
        инстанцама ($n > 1000$)?
  \item \textbf{Квантна верификација:} Развој аутоматских алата
        (на бази SMT/ATP) за верификацију исправности квантних
        алгоритама и кола.
  \item \textbf{Квантне SMT теорије:} Интеграција квантних оракла
        директно у SMT оквире као нове теорије (нпр. ``теорија
        квантне оптимизације'').
  \item \textbf{Хибридни решавачи:} Систематска истраживања
        хибридних алгоритама где Z3 (или CDCL) води квантно
        истраживање простора решења.
\end{itemize}

% TODO: Завршити закључак личним освртом на значај теме и
% потенцијалним утицајем квантних рачунара на будућност
% формалне верификације и аутоматског резоновања.
